\section{Protobuf Transformation}\label{sec:proto-transform}

At the core of the Protobuf compatibility methodology is a transformation which
projects the view of a message under one descriptor to a related message under
another descriptor.

\subsection{Origins}

Using a function to transform the value isn't an arbitrary technical choice, nor
is it only useful to model potential descriptor changes to specific Protobuf
value. It arises naturally from the proofs of the intermediate format, and that
evolution is described here.

The first and strongest intermediate parser compatibility theorem is the
schema-correct one, Theorem~\ref{thm:inter-schema-correct}, where the relations
at play are strong enough that we can directly prove that the resulting value
from a serializer / parser round-trip under the same descriptor is equal. No
transformation is needed here.

However, the next step was to weak the compatibility antecedent to arrive at
Theorem~\ref{thm:inter-id-compat} where the descriptor is still the same but can
have no relation to the message being operated on. This is no longer possible
with our new totality conditions (Section~\ref{sec:value-totality}), but led to
the first example of a transformation function. The Intermediate ID
Compatibility relation has the \textsc{M-Drop-Unknown} rule, with no constrain
which value is dropped, and \textsc{M-Missing}, which can insert a new value
into the resulting message. Because of these rules, the induction hypothesis of
the theorem by itself is too weak to prove the whole theorem since the
compatibility relation has multiple $\m_2$ related to the same $\m_1$, $\dc_1$ and
$\dc_2$ so the proof must find a witness. Instead we use the transformation to
strengthen the inductive hypothesis into naming that witness and prove that
anything from the transformation will land inside the compatibility relation.

The final intermediate format theorem, Theorem~\ref{thm:inter-compat},
reinforces this. It uses the same technique with a much more complex
transformation function that actually mirrors the Protobuf transformation. Both
are a structural recursion on the reader's descriptor field list, building the
new interpretation of the value under the new descriptor using the writer's
descriptor. The Protobuf one also has the stronger justification of enforcing
the totality requirement.

\subsection{Definition}

The full definition of the transformation function is omitted, since it is not
fully defined at the time of this writing. However, it is broken down into
levels, just like the various size functions.

\begin{minted}{lean}
Value.reinterpret        (d₁ d₂ : Desc) (v : Value) : Value
Value.reinterpretEntries (d₁ : Desc) (v : Value) :
    List ((_:Int) × Field) → List ((_:Int) × Val)
Value.reinterpretAt      (d₁ : Desc) (v : Value) (k : Int) (f₂ : Field) : Val
Slot.reinterpret         (f₁ f₂ : Field) (x : Val) : Option Val
Payload.reinterpret      (t₁ t₂ : FieldType) (p : Payload) : Option Payload
\end{minted}

The \texttt{reinterpretEntries} function is basically a \texttt{map} over the
internal association list, with the required termination metric.

The \texttt{Option} in the last two function signatures is \emph{not} a
message-level optional; returning a \texttt{none} means that the value is not
carryable by the new type and the caller falls back to the default slot value,
which is exactly what the reader sees when the writer's bytes do not populate
the field at all or using the incorrect wire type.

Mathematically, the transformation where a message $\m$ at descriptor $\dc_1$ is
transformed to $\dc_2$ is denoted $\coercion{\dc_1}{\m}{\dc_2}$.

For each key in the reading descriptor, there are three cases.

\begin{itemize}
\item \textbf{Shared.} If both descriptors carry the payload, move it over to
  the new message, transforming it if needed.
\item \textbf{Reader only.} Inject the default value for that field.
\item \textbf{Writer only.} This is a dropped field, and since we transform the
  message by walking the \emph{reader's} fields, this case cannot arise.
\end{itemize}

The other possible cases are that there is a value in the message which isn't in
the writing descriptor or that there is a field in the writing descriptor which
isn't in the message. The totality assumption from
Section~\ref{sec:value-totality} eliminates these cases from valid messages.

The recursion here is driven by the reader's descriptor since we want to produce
a message that's total over the reader's fields. Thus the termination metric is
the size that descriptor alone, not dependent on the message being
reinterpreted.

Note that the intermediate compatibility relation didn't have a drop rule in
$\dcompat$. That was because if we encountered a dropped value, we had no way to know
if the next thing in the encoding was a scalar or a nested message and the
parsing failed. Since Protobuf encoding the wire type into the value encoding,
it's always possible to skip the unknown field without de-synchronizing the byte
stream. The Protobuf drop rule will be scoped to only be applied only when
$\dom{\dc_1} \not \subset \dom{\dc_2}$.

One final note is that the measure chain used by all the transformation layers
is only \emph{weakly} decreasing within one layer, so the measures are scaled by
4 with a per-step offset to ensure that the termination metric will actually
decrease with each application. 

\subsection{Preserving \Oneof{} Groups}\label{sec:oneof-pres}

Most of the compatibility conditions in Pollux concern one field at a time. The
\Oneof{} is the exception. A reader parsing a \Oneof{} keeps only the last
member it sees, so whether a message is valid under the reader depends on
\emph{pairs} of fields. Consider a writer with two ordinary optional fields that
are both set. If the reader places those two fields in the same \Oneof{}, it
discards one of them, and the transformed message is not a valid reader message.
The writer could legitimately set both, so no field-level check can rule this
out.

The condition we need is that the reader never groups together two fields the
writer could set independently. In the writer, the only fields that can never be
set together are members of the same \Oneof{}.

\vspace{1em}
\begin{definition}[\Oneof{} Preservation]\label{def:oneof-pres}
  \[ \begin{aligned}
    \oneofpres(\dc_1, \dc_2) \triangleq{}
      & \forall\ \id_1, \id_2, \fd_1, \fd_2, \fd_1', \fd_2', g.\ \id_1 \neq \id_2 \rightarrow \\
      & \quad \dc_1[\id_1] = \fd_1 \rightarrow \dc_1[\id_2] = \fd_2 \rightarrow
              \dc_2[\id_1] = \fd_1' \rightarrow \dc_2[\id_2] = \fd_2' \rightarrow \\
      & \quad \fd_1' = \field{\oneof{g}}{\_} \rightarrow \fd_2' = \field{\oneof{g}}{\_} \rightarrow \\
      & \quad \exists\ g_0.\ \fd_1 = \field{\oneof{g_0}}{\_} \wedge \fd_2 = \field{\oneof{g_0}}{\_}
  \end{aligned} \]
\end{definition}

In words: take any two distinct field numbers that appear in both descriptors.
If the reader $\dc_2$ puts them in the same \Oneof{} $g$, then the writer
$\dc_1$ must already have put them in a single \Oneof{} $g_0$. The group
numbers $g$ and $g_0$ need not match. A group number only identifies a \Oneof{}
within its own descriptor, so a reader is free to renumber its \Oneof{}s.

The condition constrains only fields that are present in both descriptors.
This choice makes the following updates allowed.

\begin{itemize}
\item \textbf{Adding a new field to a \Oneof{}.} A reader-only field is never
  set by the writer, so it cannot conflict with anything the writer sends.
\item \textbf{Moving a single existing field into a new \Oneof{}.} Because the
  quantified field numbers must be distinct, a reader \Oneof{} with only one
  shared member is unconstrained. Its remaining members, if any, are
  reader-only. The writer sets at most that one member, so last-wins never
  discards anything. The Protobuf language guide~\cite{LanguageGuideProto}
  lists this update as safe.
\item \textbf{Splitting a \Oneof{}.} The implication runs from reader groups to
  writer groups, so a writer \Oneof{} may be divided among several reader
  \Oneof{}s, or its members may leave the \Oneof{} entirely. Either way the
  writer still sets at most one of them, which every reader grouping accepts.
\end{itemize}

The following update is ruled out.

\begin{itemize}
\item \textbf{Moving an existing field into a \Oneof{} that already holds
  another existing field}, and likewise merging two existing fields, or two
  writer \Oneof{}s, into one reader \Oneof{}. In each case some pair of shared
  fields lands in the same reader group without having been in the same writer
  group. The language guide marks these updates as unsafe, for exactly the
  reason given above.
\end{itemize}

The condition $\id_1 \neq \id_2$ is needed. Without it, instantiating both
quantified field numbers with the same $\id$ yields the requirement that every
shared field in a reader \Oneof{} was already in a writer \Oneof{}, and that
requirement forbids even the safe single-field move.

Like the rest of the compatibility relation, $\oneofpres$ looks at one layer of
a descriptor. A field whose type is a message has its own \Oneof{}s, and those
must be preserved too, because a nested message is transformed and validated in
the same way as the outer one.

\vspace{1em}
\begin{definition}[Recursive \Oneof{} Preservation]\label{def:oneof-pres-all}
  \[ \begin{aligned}
    \oneofpresall(\dc_1, \dc_2) \triangleq{}
      & \oneofpres(\dc_1, \dc_2)\ \wedge \\
      & \forall\ \id, \card_1, \dc_1', \card_2, \dc_2'.\
        \dc_1[\id] = \field{\card_1}{\dc_1'} \rightarrow
        \dc_2[\id] = \field{\card_2}{\dc_2'} \rightarrow \\
      & \oneofpresall(\dc_1', \dc_2')
  \end{aligned} \]
\end{definition}

$\oneofpresall$ requires $\oneofpres$ at the top level and at every message
field shared by both descriptors, all the way down. The recursion follows only
shared fields, because a field that exists in only one descriptor is either
skipped by the reader or initialized to its default, and neither produces a
nested message that could contain a conflict.

Theorem~\ref{thm:trans-valid} takes $\oneofpresall$ as a hypothesis, because
validity of the transformed message is checked at every depth. The compatibility
relation needs only the one-layer $\oneofpres$, as a premise of
\textsc{D-Update}. A derivation of $\dc_1 \dcompat \dc_2$ uses \textsc{D-Update}
at every pair of nested message descriptors it relates, so the derivation as a
whole establishes $\oneofpresall$ for every pair it touches.

\subsection{Theorems}

\vspace{1em}
\begin{theorem}[Transformation Specification]
  For all messages $\m$ and descriptors $\dc_1$, $\dc_2$, each value of the
  transformed message is equal to the result of transforming that value.
  \[ \coercion{\dc_1}{\m}{\dc_2}[\id] = \coercion{\dc_1}{\m[\id]}{\dc_2} \]
\end{theorem}

\vspace{1em}
\begin{theorem}[Transform WF]\label{thm:trans-wf}
  \[ \forall\ \m, \dc_1, \dc_2.\ \wfd{\dc_2} \rightarrow \wfm{\coercion{\dc_1}{\m}{\dc_2}} \]
\end{theorem}

\vspace{1em}
\begin{theorem}[Transform Total]\label{thm:trans-total}
  \[ \forall\ \m, \dc_1, \dc_2.\ \dom{\dc_2} = \dom{\coercion{\dc_1}{\m}{\dc_2}} \]
\end{theorem}

\vspace{1em}
\begin{theorem}[Transform Self]\label{thm:trans-self}
  \[ \forall\ \m, \dc.\ \wfd{\dc} \rightarrow \validv{\m}{\dc} \rightarrow \coercion{\dc}{\m}{\dc} = \m \]
\end{theorem}

\vspace{1em}
\begin{theorem}[Transform Valid]\label{thm:trans-valid}
  \[ \forall\ \m, \dc_1, \dc_2.\ \validv{\m}{\dc_1} \rightarrow \wfd{\dc_2} \rightarrow \oneofpresall(\dc_1, \dc_2) \rightarrow \validv{\coercion{\dc_1}{\m}{\dc_2}}{\dc_2} \]
\end{theorem}

The two important theorems here are Theorem~\ref{thm:trans-self}, which states
that transforming a message against it's own transform produces the same
message, and Theorem~\ref{thm:trans-valid} which we can use to show that a
transformed value is valid.

Most of the premises for these theorems should be evident, except the
\Oneof{} conditions for Theorem~\ref{thm:trans-valid}. This premise is
required, and to see why consider that a writer should legitimately set two
fields outside a \Oneof{} and then transform to a reader where those two
fields are part of the same \Oneof. The OneofPreserved function states
that any two reader keys sharing a \Oneof{} group were already in a group
together in the writer. Note that this is vacuous if either key is reader-only,
so adding a new field to an existing \Oneof{} is acceptable, while moving
an existing field into an existing \Oneof{} isn't. This is part of the
Protobuf best practices guide~\cite{LanguageGuideProto}, represented as proof
obligation here. The direction is also important, since it does allow us to
split a \Oneof{} into multiple \Oneof{} groups without
issue. Additionally, the recursive version is required to prevent the example
case from appearing in a nested message.

Now, it is possible to move an existing field into an existing \Oneof{} in
a Protobuf descriptor. In that case, you'd need to be sure that the value didn't
set multiple fields in a \Oneof{} of the reader (we'll know in advance
that it can't have multiple fields set in a writer's \Oneof). Then you
might be able to weaken the $\oneofpresall(\dc_1, \dc_2)$ to
\[ \forall\ \id_1, \id_2, g.\ \dc_2[\id_1] = \field{\oneof{g}}{\_} \rightarrow \dc_2[\id_2] = \field{\oneof{g}}{\_} \to
  \m[\id_1] = \none \vee \m[\id_2] = \none \]
Which is wouldn't be true for all messages under $\dc_1$, but might be provable
as a weakened version of Theorem~\ref{thm:trans-valid} for some systems.

{\color{red} This is another place were the functional transformation breaks
  down. In the lean serializer, everything is serialized in sorted field key
  order, so the last writer wins rule (Section~\ref{sec:last-writer-wins}) would
  be deterministic. However, since a real serializer could serialize in an
  arbitrary order, this is actually a non-deterministic operation. The
  relational encodes definition would allow this and further complicate the
  non-canonical Protobuf encoding.}

\subsection{Faithfulness Theorem}

The Protobuf compatibility framework might drop the message compatibility
relation in favor of only having the transformation. In that case, we'd end up
needing to argue that the transformation is total over all possible messages
being converted between two descriptors, i.e. $\coercion{\dc_1}{\m}{\dc_2}$ produces
some $\m'$ for all $\validv{\m}{\dc_2}$.

The current transformation function is \emph{always total} by falling back to
injecting field type defaults, so in the context of Protobuf, we want to show
that we can produce a total transformation without needing this behavior outside
of injecting the new values for new fields (since any other place were a default
value is injected would imply that something failed to reinterpret). Note that
this theorem can't be solely stated about the output of
$\coercion{\dc_1}{\m}{\dc_2}$ since the output can't differentiate between the
default value naturally occurring and being injected due to a problematic
update.

\begin{minted}{lean}
theorem Slot.reinterpret_isSome {f₁ f₂ : Field} {x : Slot}
    (h : f₁ ∝ f₂) (hm : Slot.Matches x f₁) :
    (Slot.reinterpret f₁ f₂ x).isSome
\end{minted}

Including \texttt{Slot.Matches} is important since incompatible cardinalities or
types between the slot and the original descriptor would invoke the default
value fallback anyways. This could be built into a faithfulness theorem.

\vspace{1em}
\begin{theorem}[Faithfulness]\label{thm:faith}
  For two compatible descriptors, $\dc_1 \dcompat \dc_2$ and a valid value $\m$ under
  $\dc_1$, $\validv{\m}{\dc_1}$ such that for all keys $\id$, $\dc_1[\id] = \fd_1$ and $\dc_2[\id] =
  \fd_2$ there must exist a value $x$ and $y$ such that
  \begin{itemize}
  \item $\m[\id] = x$
  \item $\coercion{\fd_1}{\m}{\fd_2} = y$
  \item $\coercion{\dc_1}{\m}{\dc_2}[\id] = y$
  \end{itemize}
\end{theorem}

Which basically characterizes the reinterpretation process as point-wise over the
message.

\begin{minted}{lean}
theorem Value.get?_reinterpret_of_compat {d₁ d₂ : Desc} {v : Value} {k f₁ f₂}
    (h : d₁ ≪ d₂) (hv : Value.Valid d₁ v)
    (h₁ : d₁.get? k = some f₁) (h₂ : d₂.get? k = some f₂) :
    ∃ x y, v.get? k = some x ∧ Slot.reinterpret f₁ f₂ x = some y ∧
      (Value.reinterpret d₁ d₂ v).get? k = some y
\end{minted}

Some corollaries from this include the fact that repeated slots don't drop
values. If the scalar transformation is only the identity function, which is
presently the case, we also have that
$\coercion{\dc_2}{\coercion{\dc_1}{\m}{\dc_2}}{\dc_{2}} = \m$ since the newly injected
fields are simply dropped in the backwards directly. This means that no data is
lost, but the statement is too strong to handle scalar conversion and in
particular integer width changes which might truncate values.

\vspace{1em}
\begin{theorem}[Point-wise Field Compatibility]\label{thm:pointwise-compat}
  \[ \wfd{\dc_1} \rightarrow \forall\ \id \in \dom{\dc_1} \cap \dom{\dc_2}. \dc_1[\id] \fieldrel \dc_2[\id] \]
\end{theorem}

Combined with Theorem~\ref{thm:faith}, Theorem~\ref{thm:pointwise-compat} states
that the descriptor relation $\dcompat$ is a maximal relation or equivalently
the weakest precondition for lossless data transmission. Again, this will be
complicated by adding integer truncation. Once that feature is added, this
concept is split into \emph{carriage}, which thinks about data being present,
and \emph{lossless-ness} that the transformed view of the data is unique.

% Local Variables:
% citar-bibliography: ("../pollux.bib")
% TeX-master: "../pollux.tex"
% TeX-command-default: "Make"
% jinx-local-words: "protobuf"
% End:
